% ============================================================
%  ch08 —— §5.3–5.5（原文 p.31–p.37）
% ============================================================
\routenote{p.31--p.37}

\subsection{拟极小 Schubert 簇}

对余权 $\lambda\in X_*(T)$，记 $P_\lambda\subset G$ 为由 $T$ 与满足 $\langle\alpha,\lambda\rangle\ge 0$ 的根空间 $U_\alpha$ 生成的抛物子群。对每个根 $\alpha$ 与 $i\in\mathbb{Z}$，记 $U_{\alpha,\ge i}\subset U_\alpha((\tau))$ 为其 $R$-点为 $U_\alpha(\tau^iR[[\tau]])$ 的子群，并令 $U_{\alpha,i}=U_{\alpha,\ge i}/U_{\alpha,\ge i+1}$。

现设 $\lambda\in X_*(T)^+$，考虑（开）Schubert 簇 $\mathrm{Gr}_\lambda\subset\mathrm{Gr}$。由 [31, Lemme 2.3 及其前的讨论]，有（固定 $\Phi^+$ 的一个全序）
\begin{equation}
\frac{G}{P_{-\lambda}}\times\prod_{\langle\alpha,\lambda\rangle\ge 2}U_{\alpha,\ge 1}\big/U_{\alpha,\ge\langle\alpha,\lambda\rangle}\xrightarrow{\ \sim\ }\mathrm{Gr}_\lambda,\qquad (g,u)\mapsto gu\tau^\lambda.
\end{equation}
$P_{-\lambda}$ 在 (5.21) 中乘积上的作用如下：先经 $P_{-\lambda}$ 在 $\prod_\alpha U_{\alpha,\ge 1}$ 上的伴随作用，再投影到满足 $\langle\alpha,\lambda\rangle\ge 2$ 的因子 $U_{\alpha,\ge 1}$。

现考虑特殊情形 $\lambda=\theta^\vee$，即支配短余根。我们需要回忆 [31] 中拟极小 Schubert 簇 $\mathrm{Gr}_{\theta^\vee}$ 的描述。记 $P_{-\theta}:=P_{-\theta^\vee}$。

\begin{lemma}
开子簇 $\mathrm{Gr}^{\mathrm{triv}}_{\theta^\vee}\subset\mathrm{Gr}_{\theta^\vee}$ 由单个 $G$-轨道组成。更精确地，态射
\[
\frac{G}{P_{-\theta}}\times U^\times_{-\theta,-1}\to G[\tau^{-1}]_1=\mathrm{Gr}^{\mathrm{triv}},\qquad (g,u_{-\theta}(c\tau^{-1}))\mapsto\mathrm{Ad}(g)u_{-\theta}(c\tau^{-1})
\]
给出到 $\mathrm{Gr}^{\mathrm{triv}}_{\theta^\vee}$ 的同构（这里 $P_{-\theta}$ 经伴随作用作用在 $U_{-\theta,-1}$ 上，而 $U^\times_{\alpha,i}=U_{\alpha,i}\setminus\{1\}$）。
\end{lemma}

\begin{proof}
这本质上是 [31, Lemme 7.2]。把 (5.21) 用于 $\lambda=\theta^\vee$，发现 (5.21) 中的乘积只剩一项 $U_{\theta,1}$（因为对任何根 $\alpha\neq\theta$ 有 $\langle\alpha,\theta^\vee\rangle\le 1$）。因此
\[
\mathrm{Gr}^{\mathrm{triv}}_{\theta^\vee}\cong\frac{G}{P_{-\theta}}\times U^{\mathrm{triv}}_{\theta,1}\cdot\tau^{\theta^\vee},
\]
其中 $U^{\mathrm{triv}}_{\theta,1}\subset U_{\theta,1}$ 是某个 $P_{-\theta}$-稳定的真开子集（$\mathrm{Gr}^{\mathrm{triv}}_{\theta^\vee}$ 不可能等于 $\mathrm{Gr}_{\theta^\vee}$，因为 $\tau^{\theta^\vee}\notin\mathrm{Gr}^{\mathrm{triv}}$）。由于 $P_{-\theta}$ 经伸缩作用在 $U_{\theta,1}$ 上，$U^{\mathrm{triv}}_{\theta,1}$ 必为 $U^\times_{\theta,1}$。故
\[
\mathrm{Gr}^{\mathrm{triv}}_{\theta^\vee}\cong\frac{G}{P_{-\theta}}\times U^\times_{\theta,1}\cdot\tau^{\theta^\vee}.
\]
另一方面，在由 $\theta$ 决定的 $SL_2$-子群中有如下计算：
\begin{equation}
\begin{pmatrix}1&c\tau\\&1\end{pmatrix}\begin{pmatrix}\tau&&\\&&\\&&\tau^{-1}\end{pmatrix}
=\begin{pmatrix}\tau&c\\&\tau^{-1}\end{pmatrix}
=\begin{pmatrix}1&c^{-1}\tau^{-1}\\&1\end{pmatrix}\begin{pmatrix}\tau&c\\&-c^{-1}\end{pmatrix}
\end{equation}
它表明：当 $c$ 可逆时，$u_\theta(c\tau)\tau^{\theta^\vee}\in\mathrm{Gr}$ 等于 $u_{-\theta}(c^{-1}\tau^{-1})\in U^\times_{-\theta,-1}\subset G[\tau^{-1}]_1$。引理得证。
\end{proof}

$G$ 的 Bruhat 分解 $G=\bigsqcup_{w\in W/W_\theta}UwP_{-\theta}$ 给出分解
\[
\frac{G}{P_{-\theta}}\times U^\times_{-\theta,-1}=\bigsqcup_{w\in W/W_\theta}\frac{UwP_{-\theta}}{P_{-\theta}}\times U^\times_{-\theta,-1}.
\]
层 $UwP_{-\theta}/P_{-\theta}\times U^\times_{-\theta,-1}$ 在 $G[\tau^{-1}]_1$ 中的像是 $\mathrm{Ad}(U)U^\times_{-w\theta,-1}$。由于 $w\mapsto-w\theta$ 给出 $W/W_\theta$ 与 $G$ 的长根集合之间的双射，可把上述分解改写为
\[
\mathrm{Gr}^{\mathrm{triv}}_{\theta^\vee}=\bigsqcup_{\beta\ \text{长根}}\mathrm{Ad}(U)U^\times_{\beta,-1}.
\]
对每个 $w\in W/W_\theta$，有同构（固定 $\Phi^+$ 的全序）
\[
\prod_{\alpha\in\Phi^+,\ \langle w^{-1}\alpha,\theta^\vee\rangle>0}U_\alpha\xrightarrow{\ \sim\ }UwP_{-\theta}/P_{-\theta}.
\]
因此，对 $\beta=-w\theta$，层 $\mathrm{Ad}(U)U^\times_{\beta,-1}$ 可写成
\begin{equation}
\begin{aligned}
\mathrm{Ad}(U)U^\times_{\beta,-1}&\xrightarrow{\ \sim\ }\prod_{\alpha\in\Phi^+,\ \langle\alpha,\beta^\vee\rangle<0}U_\alpha\times U^\times_{\beta,-1},\\
\mathrm{Ad}\Big(\prod u_\alpha\Big)u_\beta(c\tau^{-1})&\leftrightarrow\Big(\prod u_\alpha,\ u_\beta(c\tau^{-1})\Big).
\end{aligned}
\end{equation}

\begin{lemma}
引理 5.13 中的函数 $a_{-\theta}:\mathrm{Gr}^{\mathrm{triv}}_{\theta^\vee}\to\ckg_{-\theta}$ 在每个层 $\mathrm{Ad}(U)U^\times_{\beta,-1}\subset\mathrm{Gr}^{\mathrm{triv}}_{\theta^\vee}$ 上由
\[
a_{-\theta}\big(\mathrm{Ad}(u)u_\beta(c\tau^{-1})\big)=
\begin{cases}
c\,x_{-\theta},&\beta=-\theta,\\
0,&\text{否则}
\end{cases}
\]
给出，其中 $x_{-\theta}\in\ckg_{-\theta}$ 对应 $u_{-\theta}(1)\in U_{-\theta}$。特别地，
\begin{equation}
\mathrm{Gr}^{\mathrm{triv},a=0}_{\theta^\vee}=\bigsqcup_{\beta\ \text{长根},\ \beta\neq-\theta}\mathrm{Ad}(U)U^\times_{\beta,-1}.
\end{equation}
\end{lemma}

\begin{proof}
按定义，$a_{-\theta}$ 在 $\mathrm{Ad}(u)u_\beta(c\tau^{-1})$ 处的值为
\[
\frac{d}{d(\tau^{-1})}\mathrm{Ad}(u)u_\beta(c\tau^{-1})\Big|_{\tau^{-1}=0}\in\mathrm{Ad}(U)\ckg_\beta.
\]
若 $\beta\neq-\theta$，$\mathrm{Ad}(u)\ckg_\beta$ 只涉及 $\ge\beta$ 的根。若 $\beta=-\theta$，上述导数等于 $\mathrm{Ad}(u)c\,x_{-\theta}\in\ckg$，其 $\ckg_{-\theta}$-分量为 $cx_{-\theta}\in\ckg_{-\theta}$。
\end{proof}

\subsection{定理 4(1) 的证明}

把推论 5.20 用于 $V=V_{\theta^\vee}$（它只有一个非零支配权 $\theta^\vee$），得
\[
\chi_c(\PP^1\setminus\{0,\infty\},\Kl^{\theta^\vee}_{\ckG}(\varphi,\chi))=\chi_c\big(\mathrm{Gr}^{\mathrm{triv},a=0}_{\theta^\vee},\mathrm{ev}^*J\big).
\]
由引理 5.25 中的分解 (5.26)，只需对长根 $\beta\neq-\theta$ 计算 $\chi_c(\mathrm{Ad}(U)U^\times_{\beta,-1},\mathrm{ev}^*J)$。于是定理 4(1) 由下述断言推出。

\medskip\noindent\textbf{断言.} 设 $\beta\neq-\theta$ 是长根，则
\[
\chi_c\big(\mathrm{Ad}(U)U^\times_{\beta,-1},\mathrm{ev}^*J\big)=
\begin{cases}
-1,&\beta\ \text{是长单根},\\
0,&\text{否则}.
\end{cases}
\]

为证该断言，分三种情形。

\emph{情形 I：$\beta$ 为正但非单根。} 由于 $\beta>0$，有 $\mathrm{ev}(\mathrm{Ad}(U)U^\times_{\beta,-1})\subset U$。由于 $\beta$ 非单根，$\mathrm{ev}(\mathrm{Ad}(U)U^\times_{\beta,-1})$ 在 $U/[U,U]$ 中的像是平凡的，故 $\mathrm{ev}^*J$ 是 $\mathrm{Ad}(U)U^\times_{\beta,-1}$ 上的常层。由 (5.24)，$\mathrm{Ad}(U)U^\times_{\beta,-1}$ 有一个因子 $U^\times_{\beta,-1}\cong\mathbb{G}_m$，故
\[
\chi_c\big(\mathrm{Ad}(U)U^\times_{\beta,-1},\mathrm{ev}^*J\big)=\chi_c\big(\mathrm{Ad}(U)U^\times_{\beta,-1}\big)=0.
\]

\emph{情形 II：$\beta=\alpha_i$ 是长单根。} 对 $u\in U$，$\mathrm{ev}(\mathrm{Ad}(u)u_\beta(c\tau^{-1}))$ 在 $U/[U,U]$ 中的像是 $u_{\alpha_i}(c)$。按 (5.24) 的坐标，$\mathrm{ev}^*J$ 是从 $U^\times_{\beta,-1}\cong U^\times_{\alpha_i}$-因子拉回的 $\mathrm{AS}_{\varphi_i}$。故
\[
\chi_c\big(\mathrm{Ad}(U)U^\times_{\beta,-1},\mathrm{ev}^*J\big)=\chi_c\big(U^\times_{\alpha_i},\mathrm{AS}_{\varphi_i}\big)=-1.
\]

\emph{情形 III：$\beta$ 为负且 $\beta\neq-\theta$。} 由于 $\beta\neq-\theta$，存在单根 $\alpha_i$ 使 $\beta-\alpha_i$ 仍是根。则 $\alpha=-\beta+\alpha_i$ 是正根。又因 $\alpha_i\neq\beta$ 蕴涵 $\langle\alpha_i,\beta^\vee\rangle\le 1$，故
\[
\langle\alpha,\beta^\vee\rangle=\langle-\beta+\alpha_i,\beta^\vee\rangle=-2+\langle\alpha_i,\beta^\vee\rangle<0.
\]
于是 $U_\alpha$ 出现在分解 (5.24) 中。

利用 (5.24)，把 $\mathrm{Ad}(U)U^\times_{\beta,-1}$ 中的元素写成 $\mathrm{Ad}(u)u_\beta(c_\beta\tau^{-1})$，其中 $u=u_\alpha u_{\underline\alpha}$，$u_\alpha\in U_\alpha$ 而 $u_{\underline\alpha}\in A:=\prod_{\alpha'>0,\ \langle\alpha',\beta^\vee\rangle<0,\ \alpha'\neq\alpha}U_{\alpha'}$。注意
\begin{equation}
\mathrm{ev}\big(\mathrm{Ad}(u)u_\beta(c_\beta\tau^{-1})\big)=u\,u_\beta(c_\beta)\,u^{-1}=u_\alpha[u_{\underline\alpha}(c_\beta),u_\beta(c_\beta)]\,u_\beta(c_\beta)\,u_{\underline\alpha}^{-1}.
\end{equation}
由于 $\alpha\neq\pm\beta$，可应用 Chevalley 交换子关系 [9, p.36,(4)]，断言 $[u_{\underline\alpha}(c_\beta),u_\beta(c_\beta)]$ 是根群 $U_{i\underline\alpha+j\beta}$（$i,j\in\mathbb{Z}_{>0}$）中元素之积。我们的假设 (i) $\alpha+\beta$ 是单根、(ii) $\beta$ 是长根蕴涵：任何这样的 $i\underline\alpha+j\beta$（若为根）都是正根，且此形式中唯一的单根是 $\alpha+\beta$。因此 $[u_{\underline\alpha}(c_{\underline\alpha}),u_\beta(c_\beta)]\in U$，且其在 $U/[U,U]$ 中的像是 $u_{\alpha_i}(\epsilon c_{\underline\alpha}c_\beta)$，其中 $\epsilon=\pm 1$。

由 (5.27)，$\mathrm{Ad}(u)u_\beta(c_\beta)\in UB^{\mathrm{opp}}$ 当且仅当 $u_\beta(c_\beta)u_{\underline\alpha}^{-1}\in UB^{\mathrm{opp}}$。此外，当 $\mathrm{Ad}(u)u_\beta(c_\beta)\in UB^{\mathrm{opp}}$ 时，它在 $U/[U,U]$ 中的像是 $u_{\alpha_i}(\epsilon c_{\underline\alpha}c_\beta)$ 乘以另一个只依赖于 $c_\beta$ 与 $u_{\underline\alpha}$ 的元素。

在分解 (5.24) 下，$\mathrm{Ad}(U)U^\times_{\beta,-1}\cong U_\alpha\times A$。设 $\mathrm{pr}_A:\mathrm{Ad}(U)U^\times_{\beta,-1}\to A$ 为投影。由上述讨论，$\mathrm{ev}^*J$ 限制在 $\mathrm{pr}_A$ 的每个纤维上同构于 $U_\alpha$ 上的 Artin--Schreier 层。因此 $\mathrm{pr}_{A,!}\mathrm{ev}^*J=0$，从而
\[
H^*_c\big(\mathrm{Ad}(U)U^\times_{\beta,-1},\mathrm{ev}^*J\big)=H^*_c\big(A,\mathrm{pr}_{A,!}\mathrm{ev}^*J\big)=0.
\]
断言得证，定理 4(1) 的证明完成。

\subsection{伴随 Schubert 簇}

本小节设 $G$ 非单链。始终假设 $\operatorname{char}(k)$ 对 $G$ 良好，即当 $G$ 型为 $B_n,C_n$ 时 $\operatorname{char}(k)>2$，当 $G$ 型为 $F_4,G_2$ 时 $\operatorname{char}(k)>3$。

设 $\gamma$ 为 $G$ 的短支配根，$\gamma^\vee$ 为相应的长余根。完全类似 $P_{-\theta},W_\theta$ 等来定义 $P_{-\gamma},W_\gamma$ 等。按 $\langle\alpha,\gamma^\vee\rangle$（$\alpha\in\Phi$）的可能取值，有两种情形：
\begin{itemize}[leftmargin=1.6em, itemsep=.15em]
\item 型 $(B_n,C_n,F_4)$：$|\langle\alpha,\gamma^\vee\rangle|=0,1,2$；
\item 型 $(G_2)$：$|\langle\alpha,\gamma^\vee\rangle|=0,1,2,3$。
\end{itemize}
记 $\Phi^\gamma_n:=\{\alpha\in\Phi\mid\langle\alpha,\gamma^\vee\rangle=n\}$。

\begin{lemma}
设 $G$ 型为 $B_n,C_n$ 或 $F_4$，则 $\mathrm{Gr}^{\mathrm{triv}}_{\gamma^\vee}$ 由单个 $G$-轨道组成。更精确地，令
\[
V_{-\gamma}:=\prod_{\langle\alpha,\gamma^\vee\rangle=2}U_{-\alpha,-1},
\]
并把 $V_{-\gamma}$ 与其 Lie 代数等同。$P_{-\gamma}$ 在 $V_{-\gamma}$ 上的伴随作用（经其 Levi 因子 $L_\gamma$ 分解）稳定由 $V_{-\gamma}$ 上一个非退化二次型 $q_{-\gamma}$ 定义的唯一四元超曲面 $Q_{-\gamma}$，使 $L_\gamma$ 的作用经 $\mathrm{GO}(V_{-\gamma},q_{-\gamma})$ 分解\footnote{$\mathrm{GO}(V_{-\gamma},q_{-\gamma})$ 由 $V_{-\gamma}$ 的把 $q_{-\gamma}$ 保持到相差一个纯量的可逆线性自同构组成。}，并以 $V_{-\gamma}-Q_{-\gamma}$ 为稠密轨道。自然态射
\begin{equation}
\frac{G}{P_{-\gamma}}\times(V_{-\gamma}-Q_{-\gamma})\ni(g,v)\mapsto\mathrm{Ad}(g)v\in G[\tau^{-1}]_1
\end{equation}
给出同构 $\frac{G}{P_{-\gamma}}\times(V_{-\gamma}-Q_{-\gamma})\cong\mathrm{Gr}^{\mathrm{triv}}_{\gamma^\vee}$。
\end{lemma}

\begin{proof}
当 $G$ 型为 $B,C$ 或 $F_4$ 时，(5.21) 中的乘积化为 $\Lambda=\prod_{\langle\alpha,\gamma^\vee\rangle=2}U_{\alpha,1}$（一个交换幂幺群），可与其 Lie 代数等同。由于 $\mathrm{Gr}^{\mathrm{triv}}_{\gamma^\vee}$ 在 $G$-共轭下稳定，它形如 $\frac{G}{P_{-\gamma}}\times\Lambda^{\mathrm{triv}}$，其中 $\Lambda^{\mathrm{triv}}\subset\Lambda$ 是某个 $P_{-\gamma}$-稳定开子集。$P_{-\gamma}$ 在 $\Lambda$ 上的作用经 Levi 商 $L_\gamma$ 分解。

向量空间 $\Lambda$ 带双线性型
\begin{equation}
(x,y)_\Lambda:=(x,\mathrm{Ad}(w_\gamma)y)_{\ckg},
\end{equation}
其中 $w_\gamma$ 是 $\begin{pmatrix}0&1\\-1&0\end{pmatrix}$ 在对应根 $\gamma$ 的同态 $SL_2\to G$ 下的像，$(\cdot,\cdot)_{\ckg}$ 是 $\ckg$ 上 $\mathrm{Ad}(G)$-不变的非退化对称双线性配对（当 $\operatorname{char}(k)$ 良好时存在，见 [8, §1.16]）。容易验证 $(\cdot,\cdot)_\Lambda$ 是非退化对称双线性配对，且 $L_\gamma$ 在 $\Lambda$ 上的作用把该配对保持到相差纯量。记 $q_\Lambda$ 为 $(\cdot,\cdot)_\Lambda$ 相伴的二次型。

经快速的逐型验证，可知作用映射 $L_\gamma\to\mathrm{GO}(\Lambda,q_\Lambda)$ 是满的。因此 $\Lambda$ 含唯一的 $L_\gamma$-稳定不可约除子 $Q_\Lambda=\{q_\Lambda=0\}$，其补集是单个 $L_\gamma$-轨道。

另一方面，补集 $\mathrm{Gr}_{\le\gamma^\vee}-\mathrm{Gr}^{\mathrm{triv}}_{\le\gamma^\vee}$ 是代表行列式线丛类的丰沛除子，故余维为 $1$。由于 $\mathrm{Gr}_{\le\gamma^\vee}-\mathrm{Gr}_{\gamma^\vee}$ 的余维至少为 $2$，$\mathrm{Gr}_{\gamma^\vee}-\mathrm{Gr}^{\mathrm{triv}}_{\gamma^\vee}$ 在 $\mathrm{Gr}_{\gamma^\vee}$ 中余维也是 $1$。于是 $\Lambda-\Lambda^{\mathrm{triv}}$ 在 $\Lambda$ 中余维为 $1$，故必为该不可约除子 $Q_\Lambda$。这蕴涵 $\Lambda^{\mathrm{triv}}=\Lambda-Q_\Lambda$ 是单个 $P_{-\gamma}$-轨道，因此 $\mathrm{Gr}^{\mathrm{triv}}_{\gamma^\vee}$ 是单个 $G$-轨道。

与 (5.23) 相同的 $SL_2$-计算给出
\[
u_\gamma(\tau)\tau^{\gamma^\vee}=u_{-\gamma}(\tau^{-1})\in\mathrm{Gr}^{\mathrm{triv}}.
\]
因此 $\Lambda^{\mathrm{triv}}\tau^{\gamma^\vee}=\mathrm{Ad}(P_{-\gamma})u_\gamma(\tau)\tau^{\gamma^\vee}=\mathrm{Ad}(P_{-\gamma})u_{-\gamma}(\tau^{-1})\subset\mathrm{Gr}^{\mathrm{triv}}$。与上述类似的论证表明 $\mathrm{Ad}(P_{-\gamma})u_{-\gamma}(\tau^{-1})$ 是 $V_{-\gamma}=\prod_{\langle\alpha,\gamma^\vee\rangle=2}U_{-\alpha,-1}$ 中由某四元超曲面 $Q_{-\gamma}$ 的补集所定义的开子集。因此
\[
\mathrm{Gr}^{\mathrm{triv}}_{\gamma^\vee}=\frac{G}{P_{-\gamma}}\times(\Lambda-Q_\Lambda)\tau^{\gamma^\vee}=\frac{G}{P_{-\gamma}}\times(V_{-\gamma}-Q_{-\gamma})\subset G[\tau^{-1}]_1.
\]
\end{proof}

\begin{lemma}
态射 (5.29) 扩充为消解
\[
\nu:\frac{G}{P_{-\gamma}}\times V_{-\gamma}\to\mathrm{Gr}^{\mathrm{triv}}_{\le\gamma^\vee},
\]
它由引理 5.28 在 $\mathrm{Gr}^{\mathrm{triv}}_{\gamma^\vee}$ 上是同构。且
\begin{enumerate}[label=(\arabic*), leftmargin=1.9em, itemsep=.2em]
\item 纤维 $\nu^{-1}(\star)\cong G/P_{-\gamma}$；
\item 在 $\mathrm{Gr}^{\mathrm{triv}}_{\theta^\vee}$ 上的纤维同构于 $L_\theta/L_\theta\cap P_{-\gamma}$。
\end{enumerate}
\end{lemma}

\begin{proof}
由于 $\frac{G}{P_{-\gamma}}\times(V_{-\gamma}-Q_{-\gamma})$ 在 $\frac{G}{P_{-\gamma}}\times V_{-\gamma}$ 中稠密，而 $\mathrm{Gr}^{\mathrm{triv}}_{\le\gamma^\vee}$ 是 $\mathrm{Gr}^{\mathrm{triv}}_{\gamma^\vee}$ 在 $\mathrm{Gr}^{\mathrm{triv}}$ 中的闭包，态射 (5.29) 扩充为 $\nu$。由引理 5.28，$V_{-\gamma}$ 上有三个 $\mathrm{Ad}(P_{-\gamma})$（或 $\mathrm{Ad}(L_\gamma)$）-轨道：$V_{-\gamma}-Q_{-\gamma}$、$Q^\times_{-\gamma}:=Q_{-\gamma}-\{0\}$ 与 $\{0\}$，它们给出 $\frac{G}{P_{-\gamma}}\times V_{-\gamma}$ 的三个 $G$-轨道。

轨道 $\frac{G}{P_{-\gamma}}\times\{0\}$ 映到 $\mathrm{Gr}^{\mathrm{triv}}_0=\star$，这证明 (1)。

于是轨道 $\frac{G}{P_{-\gamma}}\times Q^\times_{-\gamma}$ 必映到 $\mathrm{Gr}^{\mathrm{triv}}_{\theta^\vee}$。$u_{-\theta}(\tau^{-1})$ 在 $\frac{G}{P_{-\gamma}}\times Q^\times_{-\gamma}$ 中与在 $\mathrm{Gr}^{\mathrm{triv}}_{\theta^\vee}=\frac{G}{P_{-\theta}}\times U^\times_{-\theta,-1}$ 中的 $G$-稳定子分别是 $P_{-\gamma}\cap P^1_{-\theta}$ 与 $P^1_{-\theta}$，其中 $P^1_{-\gamma}=\ker(\theta:P_{-\theta}\to\mathbb{G}_m)$。故纤维 $\nu^{-1}(u_{-\theta}(\tau^{-1}))$ 等于 $P^1_{-\theta}/P^1_{-\theta}\cap P_{-\gamma}$。容易验证包含 $L_\theta\hookrightarrow P_{-\theta}$ 与 $P^1_{-\theta}\hookrightarrow P_{-\theta}$ 诱导同构
\[
L_\theta/L_\theta\cap P_{-\gamma}\xrightarrow{\ \sim\ }P_{-\theta}/P_{-\theta}\cap P_{-\gamma}\xleftarrow{\ \sim\ }P^1_{-\theta}/P^1_{-\theta}\cap P_{-\gamma}.
\]
由引理 5.22，$\mathrm{Gr}^{\mathrm{triv}}_{\theta^\vee}$ 是单个 $G$-轨道，故 $\nu$ 在 $\mathrm{Gr}^{\mathrm{triv}}_{\theta^\vee}$ 上的所有纤维都同构于 $\nu^{-1}(u_{-\theta}(\tau^{-1}))\cong L_\theta/L_\theta\cap P_{-\gamma}$。
\end{proof}

$G$ 的 Bruhat 分解 $G=\bigsqcup_{W/W_\gamma}UwP_{-\gamma}$ 给出分解
\[
\frac{G}{P_{-\gamma}}\times V_{-\gamma}=\bigsqcup_{W/W_\gamma}\frac{UwP_{-\gamma}}{P_{-\gamma}}\times V_{-\gamma}.
\]
映射 $w\mapsto-w\gamma$ 给出 $W/W_\gamma$ 与 $G$ 的短根集合之间的双射。对 $\beta=-w\gamma$，令 $V_\beta=\mathrm{Ad}(w)V_{-\gamma}$。则上述分解可改写为
\begin{equation}
\frac{G}{P_{-\gamma}}\times V_{-\gamma}=\bigsqcup_{\beta\ \text{短根}}\mathrm{Ad}(U)V_\beta.
\end{equation}
与拟极小情形一样，可进一步把每个层写成
\begin{equation}
\begin{aligned}
\mathrm{Ad}(U)V_\beta&\xrightarrow{\ \sim\ }\prod_{\alpha\in\Phi^+,\ \langle\alpha,\beta^\vee\rangle<0}U_\alpha\times\prod_{\langle\beta',\beta^\vee\rangle=2}U_{\beta',-1},\\
\mathrm{Ad}\Big(\prod u_\alpha\Big)\prod u_{\beta'}(c_{\beta'}\tau^{-1})&\leftrightarrow\Big(\prod u_\alpha,\ \prod u_{\beta'}(c_{\beta'}\tau^{-1})\Big).
\end{aligned}
\end{equation}

\begin{remark}
\begin{enumerate}[label=(\arabic*), leftmargin=1.9em, itemsep=.2em]
\item 对每个短根 $\beta$，集合 $\Phi^\beta_2$（即出现在 $V_\beta$ 各因子中的那些根）按高度全序。为此只需证明：对不同的 $\beta',\beta''\in\Phi^\beta_2$，有 $\langle\rho^\vee,\beta'-\beta''\rangle\neq 0$。在引理 5.28 的证明中我们已注意到 $L_\gamma\twoheadrightarrow\mathrm{GO}(\Lambda,q_\Lambda)$ 是满的，这意味着对任意两个不同的根 $\alpha',\alpha''\in\Phi^\gamma_2$（即出现在 $\Lambda$ 各因子中的根），差 $\alpha'-\alpha''$ 是 $L_\gamma$ 的某个根的非零倍数。由于当 $\beta=w\gamma$ 时 $\Phi^\beta_2=w\Phi^\gamma_2$，差 $\beta'-\beta''$ 也是 $L_\beta$ 的某个根的非零倍数，故 $\langle\rho^\vee,\beta'-\beta''\rangle\neq 0$。

集合 $\Phi^\beta_2$ 带对合 $\beta'\mapsto 2\beta-\beta'$，其唯一不动点是 $\beta$。该对合反序。
\item 仍由引理 5.28 证明中的 $L_\gamma\twoheadrightarrow\mathrm{GO}(\Lambda,q_\Lambda)$，所有根 $\gamma'\in\Phi^\gamma_2\setminus\{\gamma\}$ 都被 $W_\gamma$ 置换。把 $\gamma$ 换成任一短根 $\beta$ 也有类似结论。特别地，$\Phi^\beta_2$ 中除 $\beta$ 自身外都是长根。
\end{enumerate}
\end{remark}

\begin{lemma}
对短根 $\beta$、$\prod u_{\beta'}(c_{\beta'}\tau^{-1})\in V_\beta$（乘积跑遍 $\Phi^\beta_2$）与 $u\in U$，有
\[
a_{-\theta}\Big(\mathrm{Ad}(u)\big(\prod u_{\beta'}(c_{\beta'}\tau^{-1})\big)\Big)=
\begin{cases}
c_{-\theta}\,x_{-\theta},&-\theta\in\Phi^\beta_2,\\
0,&\text{否则}.
\end{cases}
\]
\end{lemma}

\begin{proof}
证明与引理 5.25 类似。
\end{proof}
